\section{Discussion}
\label{sec:discussion}

\subsection{Key Findings and Their Interpretation}

Our confirmed results from h-e1 establish three findings with direct implications for the entropy-guided selective SWA framework.

\textbf{Finding 1: Attention concentration is a near-universal structural property of Llama-2-7B within the WikiText-103 evaluation domain.} The 100\% satisfaction rate (Gini $> 0.5$, top-10\% share $> 0.5$ across all 200 evaluation examples) suggests that heavy-hitter attention concentration is a consistent structural property of Llama-2-7B's trained weights, not an input-specific phenomenon within this evaluation domain. This has practical implications: practitioners can apply the entropy scoring criterion on any representative calibration set from the same domain and expect to observe meaningful concentration structure.

\textit{Broader implication:} Our in-domain results are consistent with the possibility that attention concentration may be architecturally grounded rather than purely input-driven, but cross-domain validation is required to confirm this. Cross-domain calibration stability remains an open question (see Limitations L2 and L5), and this generalization should not be assumed without empirical verification.

\textbf{Finding 2: Head-mean pooling is a first-class methodological choice.} The 32\% relative gap between head-mean Gini (0.681) and head-max Gini (0.466) is a methodological finding with consequences beyond this paper. Any attention analysis method that uses per-layer entropy as a signal---for pruning, routing, quantization, or structural modification---must explicitly justify its aggregation choice. Using head-max instead of head-mean risks systematically underestimating concentration and misidentifying layers as ``not locally-biased'' when they are.

\textit{Broader implication:} For layer-level characterization of multi-head attention, head-mean aggregation captures the layer's typical behavior across all heads, while head-max captures the outlier head in each layer. These are different quantities with different stabilities and different information content. The field should be explicit about which is intended.

\textbf{Finding 3: Entropy ranking stability ($\rho \geq 0.8$) enables deterministic layer selection.} The high Spearman correlation across calibration subsets means that the top-$k$ layer selection is not sensitive to the specific 100 sequences used for calibration. In practice, a practitioner can run the entropy scoring once on any available calibration data and trust that the selected layers would be the same with any other representative sample from the same domain.

\subsection{Open Questions and Motivation for Pending Experiments}

\textbf{Does the entropy criterion's structural validity translate to accuracy preservation?} H-e1 confirms that high-entropy layers have diffuse attention patterns---but whether \textit{constraining} those layers to a local window (SWA) causes accuracy degradation depends on whether those layers' diffuse patterns are globally necessary or incidentally broad. H-e2 directly tests this by converting the 4 highest-entropy layers and measuring WikiText-103 perplexity.

\textbf{Does entropy-guided selection produce lower degradation than uninformed selection?} The exploratory QA F1 probe (Section~\ref{sec:results}) shows a directional but non-significant advantage for entropy selection (0.43 pp vs.\ 0.67 pp, $p=0.4507$, $N=200$). H-m1 will provide the proper comparison with SWA masking and perplexity metric at adequate statistical power. No superiority claim is made at this stage.

\textbf{Where is the $k^*$ boundary?} H-m2 tests $k = 8$ conversion. If $k = 4$ preserves perplexity within 2 points and $k = 8$ exceeds the threshold, the boundary characterizes how many layers can be converted before accuracy degrades. This $k^*$ characterization is publishable regardless of whether $k = 4$ passes.

\subsection{Limitations}

We state the following limitations explicitly, along with their impact on paper claims.

\textbf{L1: Core accuracy-preservation claims (P1, P2, P3) are untested.} H-e2, h-m1, and h-m2 have not been executed. All accuracy-preservation claims are explicitly conditional or marked as pending. This is the primary limitation of this interim synthesis. The paper's contribution is the entropy criterion characterization (h-e1), which is meaningful independently of h-e2 outcomes.

\textit{Why acceptable:} The pipeline invoked Phase 4.5/6 before h-e2 execution. The h-e1 findings provide a meaningful checkpoint contribution. Both positive (P1 passes) and negative (P1 fails) outcomes for h-e2 are publishable.

\textbf{L2: Single model and domain scope.} All confirmed findings apply to Llama-2-7B evaluated on WikiText-103. Generalizability to Llama-3-8B (which uses grouped-query attention, requiring adapted entropy computation), encoder-decoder models, or other calibration domains is untested.

\textit{Why acceptable:} Single-model existence PoC is the stated scope of h-e1. Cross-model generalization is identified as future work.

\textbf{L3: Preliminary P2 evidence uses proxy operation and metric.} The QA F1 comparison in h-e1 uses top-$k$ attention score retention (a different operation from SWA masking) and F1 (a different metric from perplexity). The directional trend cannot be cited as evidence for P2.

\textit{Why acceptable:} Directional consistency is reported for transparency. H-m1, with the correct operation (SWA masking) and metric (perplexity), will provide definitive evidence.

\textbf{L4: Residual stream compensation mechanism is hypothesized, not empirically confirmed.} Steps 2 and 3 of the causal chain---that SWA layers do not lose needed functionality, and that 28 remaining full-attention layers compensate via residual stream---are theoretical motivations. Note that SWARR \cite{Liu2026Architecture} demonstrates that SFT alone is insufficient after SWA conversion---underscoring the challenge of fine-tuning-free conversion but not directly speaking to the selective (4-of-32 layers) setting.

\textit{Why acceptable:} Mechanism motivation is theory-grounded; h-e2 will provide the empirical test ($\Delta\text{perplexity} \leq 2.0$ is the mechanistic falsifier).

\textbf{L5: Exact Spearman $\rho$ per-pair values are not tabulated in this interim report.} The gate criterion ($\min \rho \geq 0.8$) is confirmed from h-e1 implementation, but exact per-pair point estimates were not separately tabulated. Reviewers cannot distinguish whether the stability is narrowly above threshold ($\rho \approx 0.80$--$0.82$) or substantially above it ($\rho \approx 0.93$--$0.96$). Final submission will include exact per-pair $\rho$ values extracted from h-e1 implementation logs.

\subsection{Broader Impact}

This work targets inference efficiency for large language models without fine-tuning. Positive impacts include: reducing compute cost for LLM deployment (lower energy, lower cost, broader access), enabling practitioners to apply SWA conversion on existing models without retraining infrastructure. The method requires only 100 calibration sequences and a single GPU.

Potential negative impacts are limited: the method modifies model inference behavior without changing training data or outputs in ways that could introduce harmful biases. However, practitioners should validate that SWA conversion does not disproportionately affect performance on minority-group or low-resource language examples before deploying in sensitive settings.
